\subsection{The complete leading layer in degree fourteen}
\label{subsec:fr-fourteen-normalization}

Use the pure integral $L_{11}$ of
Lemma~\ref{lem:fr-pure-eleven-certificate}, normalized by
$[S^{10}T]L_{11}=8448$ and $[S^6TST^3]L_{11}=0$.
We continue to write $E_n=C_\infty\cap F^n$ for the dyadic
word filtration.

\begin{theorem}\label{thm:fr-fourteen-normalization}
On $E_{14}$ the coordinate map
\[
 b_{14}(g)=\bigl(
 [S^8TST^4]f_g,\ 
 [S^7TST^5]f_g,\ 
 [S^6TST^2ST^3]f_g
 \bigr)
\]
is continuous and additive with scalar conjugation weight
$\lambda_2^{14}$, and
\[
 b_{14}(E_{14})=b_{14}(E_{14}\cap H)=\mathbf Z_2^3.
\]
Its kernel is $E_{15}=C_\infty\cap F^{15}$. Consequently
\[
 E_{14}/E_{15}\simeq\mathbf Z_2^3,\qquad
 G=E_{15}H,\qquad
 B\simeq E_{15}/\overline{[G,E_{15}](E_{15}\cap H)}.
\]
The original $X_2=H^1(N,\mathbf F_2)^G$ restriction kernel
to $E_{15}$ remains the odd-prime rational Kummer span plus
$\mathbf F_2\eta_3$. No nonzero linear combination of the three
coordinates modulo two on $E_{14}$ extends invariantly to $N$.
\end{theorem}

\begin{proof}
First construct the required primitive arithmetic pure
degree-eleven input without dividing a group element.
Starting with the arithmetic $u_9$ constructed in degree
twelve, form $\Delta_9(u_9)$. It starts in degree ten.
Use a profinite power of the beta-trivial arithmetic $w_{10}$
to remove its degree-ten term, obtaining a word starting in
degree eleven. Then use a profinite power of $w_{11}$ to make
its coefficient $b_{11}=[S^6TST^3]$ zero. Denote the result
by $u_{11}\in H\cap F^{11}$. The two corrections preserve
every beta moment and every rational Kummer coordinate.

The lower positive even moments already vanish. The exact
weight laws therefore give
\[
 m_{11}(u_{11})=
 \frac{(3^{11}-3^9)(3^{11}-3^7)(3^{11}-3^5)(3^{11}-3^3)}
             {1814400}\,\tau_{11}(\sigma)
 =475764757947648\,\tau_{11}(\sigma).
\]
The finite coefficient certificate now identifies its entire
leading term:
\begin{equation}\label{eq:fr-fourteen-pure-eleven-input}
 \operatorname{in}_{11}(u_{11})=\epsilon_{11}L_{11},
 \qquad
 \epsilon_{11}=\frac{619485361911}{11}\tau_{11}(\sigma)
                  \in\mathbf Z_2^\times.
\end{equation}
The denominator eleven is a unit in $\mathbf Z_2$; the only
operations on group elements were the stated profinite powers.

Define three degree-fourteen integer Lie polynomials
\[
 V_1=\{P_3,L_{11}\},\qquad
 V_2=\{L_5,L_9\},\qquad
 V_3=\{P_3,\{P_3,L_8\}\}.
\]
At the three words
\[
 \omega_1=S^8TST^4,\qquad
 \omega_2=S^7TST^5,\qquad
 \omega_3=S^6TST^2ST^3
\]
their coefficient matrix is
\begin{equation}\label{eq:fr-fourteen-unit-minor}
 \bigl([\omega_i]V_j\bigr)_{i,j}
 =
 \begin{pmatrix}
 -563792&3526&15\\
 1143317&494&-24\\
 -4406053&62987&71
 \end{pmatrix},
 \qquad \det=327441420801\in\mathbf Z_2^\times.
\end{equation}
The following table supplies the four Ihara contributions
$AB,-BA,\partial_A B,-\partial_B A$ to every entry:
\[
\begin{array}{c|r|r|r|r}
(A,B;\omega)&AB&-BA&\partial_A B&-\partial_B A\\\hline
(P_3,L_{11};\omega_1)&0&0&-760320&196528\\
(P_3,L_{11};\omega_2)&0&0&1216512&-73195\\
(P_3,L_{11};\omega_3)&0&0&-3419688&-986365\\
(L_5,L_9;\omega_1)&0&288&-38688&41926\\
(L_5,L_9;\omega_2)&0&0&57696&-57202\\
(L_5,L_9;\omega_3)&0&-21312&4003&80296\\
(P_3,\{P_3,L_8\};\omega_1)&0&0&0&15\\
(P_3,\{P_3,L_8\};\omega_2)&0&0&0&-24\\
(P_3,\{P_3,L_8\};\omega_3)&0&0&28&43
\end{array}
\]
The explicit Lyndon vectors and derivation coefficient rule
already printed, together with the complete finite certificate
for $L_{11}$, specify this calculation exactly.

The three corresponding genuine arithmetic words are
\[
 z_1=[\sigma,u_{11}],\qquad
 z_2=[u_5,u_9],\qquad
 z_3=[\sigma,w_{11}].
\]
They belong to $E_{14}\cap H$: their leading degrees are at
least fourteen, and commutators in $M$ have trivial beta and
rational Kummer coordinates. Their leading terms are
\[
 \operatorname{in}_{14}(z_1)=a\epsilon_{11}V_1,\qquad
 \operatorname{in}_{14}(z_2)=20726199\,\tau_5\tau_9V_2,
 \qquad
 \operatorname{in}_{14}(z_3)=9a^3\tau_5V_3.
\]
Here $a=m_3(\sigma)$ and $\tau_j=\tau_j(\sigma)$; every
displayed scalar is odd. The second scalar is
$9\cdot2302911$, and the last uses
$\operatorname{in}_{11}(w_{11})=9a^2\tau_5\{P_3,L_8\}$.

By Lemma~\ref{lem:full-gt-leading-term} and
\cite[Theorem~1.1, p.~2]{Willwacher}, the full rational leading
space in degree fourteen has dimension at most $d_{14}=3$.
Indeed its free Lie types are $(3,11)$, $(5,9)$, and
$(3,3,3,5)$; the apparent type $(7,7)$ vanishes. The last
listed type has one-dimensional span, represented by the
triple adjoint of the weight-three generator on the
weight-five generator.
The three genuine arithmetic leading vectors above are
independent by \eqref{eq:fr-fourteen-unit-minor}. They
therefore span the entire rational leading space.

For an integral leading polynomial in that rational space,
its three selected Magnus coefficients are integral.
The inverse of the matrix
\eqref{eq:fr-fourteen-unit-minor} is over $\mathbf Z_2$,
so its coefficients on $V_1,V_2,V_3$ all belong to
$\mathbf Z_2$. Conversely the three arithmetic words, with
their odd scalar factors, realize every integral triple
by profinite powers and products. This proves exact
integral saturation of both images, not merely equality
after tensoring with $\mathbf Q_2$.

Leading terms on $F^{14}$ add and transform by
$\lambda_2^{14}$. The three selected coefficients vanish
if and only if the entire degree-fourteen term vanishes.
Thus the kernel is $E_{15}$ and the coordinate map has
the stated properties. Arithmetic correction gives
$E_{14}=E_{15}(E_{14}\cap H)$, hence $G=E_{15}H$ and
the unchanged canonical relative quotient.

Finally all three correcting words are killed by each
$\chi\in X_2$. For the first, put
\[
 c_{3,11}=(t-t^9)(t-t^7)(t-t^5)(t-t^3)
                    \in\widehat{\mathbf Z}.
\]
The full profinite weight-one Kummer laws and the
beta-trivial corrections give
$\rho_q(u_{11})=c_{3,11}\rho_q(\sigma)$ for every
rational prime, including two. Hence
$u_{11}\sigma^{-c_{3,11}}\in D$ and
\[
 z_1=[\sigma,u_{11}\sigma^{-c_{3,11}}].
\]
For the second put $c_{5,9}=(t-t^7)(t-t^5)$.
Then $\rho_q(u_9)=c_{5,9}\rho_q(u_5)$ for all such
coordinates, so $u_9u_5^{-c_{5,9}}\in D$ and
\[
 z_2=[u_5,u_9u_5^{-c_{5,9}}].
\]
For the third, $w_{11}\in N$ already. Each is consequently
a commutator of an element of $G$ with an element of $N$.
Invariance gives $\chi(z_i)=0$ for all three. Arithmetic
correction preserves the restriction kernel from
$E_{14}$ to $E_{15}$, yielding the displayed original
kernel and excluding all nonzero invariant extensions
of the three quotient characters.
\end{proof}
